% Session 10 -- neural network quantum states: what the network actually computes, how the
% variational energy is estimated by sampling, and where NNQS sits among the ansaetze.
% Placement: sec. "Neural network quantum states", at the paragraph describing the slides'
% nested picture of the variational families.
\begin{aqmfigure}{fig:s10:nnqs}{Neural network quantum states. \textbf{(a)} What the ansatz is. The
exponentially long table of coefficients in (10.2) is not compressed, it is never written down at
all. The configuration $s$ is given as an $L$-pixel ``image,'' and the network returns the one
coefficient $C_\theta(s)$ that we ask for. This is the same structure as in
figure~\ref{fig:s10:template}, where only the input and the output have a new label. \textbf{(b)}
Training. The variational energy (10.3) is the loss, and we estimate it by sampling configurations
from the Born distribution and averaging the local energy (10.4). This is what makes the method
possible, since $E_{\rm loc}(s)$ only needs the few matrix elements $H_{ss'}$ that a sparse
Hamiltonian connects to $s$, and never the full $2^L\times2^L$ matrix. \textbf{(c)} The nested
regions drawn on the slides. A matrix product state is a highly structured special case of a PEPS,
a PEPS is in turn a state that a large enough NNQS can represent, and NNQS is the outer region
because it imposes the least structure of all on how $C$ depends on $s$. This picture organizes
the lecture and is not a theorem, since the outer inclusion is a statement about networks that are
expressive enough and not about every architecture. Note also that expressiveness has a cost,
because the order is the opposite one for how easy each family is to optimize.}

\begin{tikzpicture}[
  unit/.style={aqm, circle, draw=black, fill=white, minimum size=5.5pt, inner sep=0pt},
  unitin/.style={unit, fill=aqmfill, draw=aqmaccent},
  unitout/.style={unit, fill=aqmfill, draw=aqmaccent, minimum size=7pt},
  wire/.style={draw=aqmmuted, line width=0.22pt},
  stepbox/.style={aqmbox, text width=3.25cm, font=\footnotesize, minimum height=2.15cm},
  ring/.style={aqm, draw=black, fill=white},
]

% ===================== (a) the table the network replaces ====================
\node[aqmpanel, anchor=south west] at (0,4.55) {(a) the ansatz};

\node[aqmlabsmall, anchor=north west, align=left] (tab) at (0.42,4.35)
  {$\uparrow\uparrow\uparrow\uparrow\uparrow \;\to\; C(s)$\\[1pt]
   $\uparrow\uparrow\uparrow\uparrow\downarrow \;\to\; C(s)$\\[1pt]
   $\uparrow\uparrow\uparrow\downarrow\uparrow \;\to\; C(s)$\\[1pt]
   $\uparrow\uparrow\downarrow\uparrow\uparrow \;\to\; C(s)$\\[1pt]
   $\vdots$\\[1pt]
   $\downarrow\downarrow\downarrow\downarrow\downarrow \;\to\; C(s)$};
\draw[aqmbrace, decoration={mirror}] ([xshift=-4pt]tab.north west) -- ([xshift=-4pt]tab.south west);
\node[aqmlabsmall, anchor=north west, text width=3.0cm, align=left, text=aqmmuted]
  at ([xshift=-14pt,yshift=-4pt]tab.south west)
  {$2^L$ numbers, one per configuration: hopeless past $L\approx40$};

\draw[aqmarrowaccent, line width=1.4pt] (3.15,3.25) -- (4.05,3.25);
\node[aqmlabsmall, anchor=south, align=center, text width=1.5cm, text=aqmaccent] at (3.60,3.34)
   {not stored,\\[-2pt] generated};

% --- the network -------------------------------------------------------------
\foreach \j/\s in {0/\downarrow,1/\uparrow,2/\uparrow,3/\downarrow,4/\uparrow} {
  \node[unitin] (I\j) at (4.90,0.40*\j+2.05) {};
  \node[aqmlabsmall, anchor=east, inner sep=2pt] at (4.77,0.40*\j+2.05) {$\s$};
}
\foreach \j in {0,...,5} \node[unit] (A\j) at (5.95,0.40*\j+1.85) {};
\foreach \j in {0,...,5} \node[unit] (B\j) at (7.00,0.40*\j+1.85) {};
\node[unitout] (O0) at (8.05,3.05) {};
\node[unitout] (O1) at (8.05,2.45) {};

\begin{scope}[on background layer]
  \foreach \a in {0,...,4} \foreach \b in {0,...,5} \draw[wire] (I\a) -- (A\b);
  \foreach \a in {0,...,5} \foreach \b in {0,...,5} \draw[wire] (A\a) -- (B\b);
  \foreach \a in {0,...,5} \foreach \b in {0,1}     \draw[wire] (B\a) -- (O\b);
\end{scope}

\node[aqmlabsmall, anchor=west] at (8.25,3.05) {$\mathrm{Re}\,C_\theta$};
\node[aqmlabsmall, anchor=west] at (8.25,2.45) {$\mathrm{Im}\,C_\theta$};
\node[aqmlabsmall, anchor=north, align=center] at (6.45,1.75)
   {input $s=(s_1,\dots,s_L)$: one node per spin};

% ===================== (c) the variational hierarchy =========================
\node[aqmpanel, anchor=south west] at (10.15,4.55) {(c) where NNQS sits};

\def\hx{12.60}
\draw[ring, line width=\aqmthinwidth] (\hx,2.05) ellipse (2.55 and 2.15);
\draw[ring, draw=aqmaccent, fill=aqmfill, line width=1.2pt] (\hx,1.80) ellipse (2.05 and 1.70);
\draw[ring] (\hx,1.55) ellipse (1.60 and 1.28);
\draw[ring] (\hx,1.32) ellipse (1.15 and 0.88);
\draw[ring, dashed, line width=\aqmthinwidth] (\hx,1.12) ellipse (0.78 and 0.50);

\node[aqmlabsmall, anchor=center] at (\hx,3.78) {full Hilbert space};
\node[aqmlabsmall, anchor=center, text=aqmaccent] at (\hx,3.16) {\textbf{NNQS}};
\node[aqmlabsmall, anchor=center] at (\hx,2.55) {PEPS};
\node[aqmlabsmall, anchor=center] at (\hx,1.98) {MPS};
\node[font=\tiny, anchor=center] at (\hx,1.12) {gapped 1D};

\node[aqmlabsmall, anchor=north, align=center, text width=5.0cm, text=aqmmuted] at (\hx,-0.25)
   {more expressive outwards; easier to optimize inwards};

% ===================== (b) the sampled variational loop ======================
\node[aqmpanel, anchor=south west] at (0,-1.35) {(b) training: the energy is the loss};

\node[stepbox, anchor=north west] (s1) at (0,-1.60)
   {sample configurations from the Born rule,\\[1pt]
    $p(s)=\dfrac{|C_\theta(s)|^2}{\sum_{s'}|C_\theta(s')|^2}$};
\node[stepbox, anchor=north west] (s2) at (3.85,-1.60)
   {$E_{\rm loc}(s)=\displaystyle\sum_{s'}H_{ss'}\frac{C_\theta(s')}{C_\theta(s)}$\\[1pt]
    only $s'$ that $H$ connects to $s$};
\node[stepbox, anchor=north west] (s3) at (7.70,-1.60)
   {$E[\theta]=\big\langle E_{\rm loc}(s)\big\rangle_{s\sim p}$\\[1pt] (10.4): this is\\ the loss function};
\node[stepbox, anchor=north west, draw=aqmaccent, fill=aqmfill] (s4) at (11.55,-1.60)
   {gradient (10.5), then\\[1pt] $\theta_k\to\theta_k-\eta\,\partial E/\partial\theta_k$};

\draw[aqmarrow] (s1.east) -- (s2.west);
\draw[aqmarrow] (s2.east) -- (s3.west);
\draw[aqmarrow] (s3.east) -- (s4.west);
\draw[aqmarrowaccent] (s4.south) -- ++(0,-0.55) -| (s1.south);
\node[aqmlabsmall, anchor=south, text=aqmaccent, fill=white, inner sep=2.5pt] at (7.30,-4.30)
   {repeat until the variational energy stops falling};

\end{tikzpicture}

\end{aqmfigure}
